% Abstract (transferred from docs/paper/paper_draft.md)
\begin{abstract}
Biological connectomes are increasingly used as structural priors for artificial neural networks, but it is rarely tested whether the \emph{specific} biological wiring matters more than its coarse graph statistics. We asked whether a network constrained by real \emph{Drosophila} visual-system connectivity behaves differently from matched non-biological topologies on a human brain MRI task.

From the MaleCNS v1.0 connectome we extracted a 128-neuron, 989-edge local retinotopic patch of annotated columnar visual neurons. We used it as a fixed, sparse propagation mask between a small CNN encoder and a linear classifier. The same architecture was trained with a degree-preserving rewired version of the graph (the primary null model) and with Watts–Strogatz-type, Erdős–Rényi and Barabási–Albert controls. The task was glioma-versus-meningioma classification on the Cheng et al. Figshare dataset (171 patients, 2,134 T1-weighted contrast-enhanced slices).

The main evaluation was leakage-safe cross-validation using the official Figshare outer test folds, with 5 folds × 5 seeds per model and every patient tested exactly once. Across 25 paired fold×seed runs, the patient-level ROC AUC difference between the Drosophila graph and its rewired control was −0.002 ± 0.030 (Drosophila higher in 12 of 25 runs). At fold level it was −0.002 ± 0.008. Graph-constrained models, whether biological or not, reached higher ROC AUC than a simple CNN baseline; the biological wiring did not.

Under the evaluated setting, the Drosophila visual connectome did not provide a measurable topology-specific advantage. No statistical significance or clinical validity is claimed.

\end{abstract}
